%==============================================================================
%  PAPER 2 -- Task 2, Evolutionary and Fuzzy Systems
%
%  \input{} by main.tex (combined document) and by task2_standalone.tex.
%  Section-level content only: no preamble, no title block, no bibliography.
%==============================================================================

\section{Introduction}
\label{sec:t2-intro}

Assistive-care housing exists so that residents whose disability limits
mobility can live independently. Environmental control is central to that:
someone who cannot readily cross a room to reach a thermostat or a light
switch, or who cannot quickly add or remove a layer of clothing, depends on the
room getting it right without being asked. This paper designs a fuzzy logic
controller for the two environmental services where that dependence is
sharpest --- thermal comfort and lighting --- tunes its membership functions with
a genetic algorithm, and benchmarks three optimisation techniques against each
other on standard test functions.

Fuzzy control suits the problem for a specific reason. Zadeh's fuzzy
sets~\cite{zadeh1965} let a controller be specified in the vocabulary a domain
expert already uses --- ``the room is cool'', ``the resident is resting'' ---
rather than in a transfer function. Where the specification originates as human
judgement, and where somebody other than an engineer must approve the resulting
behaviour, that is an engineering advantage rather than a presentational one.

\subsection{Scope and contribution}

The controlled space is one room of an assistive-care flat. Four sensor inputs
drive two actuators through 54 rules. The contribution is not a new inference
algorithm; it is a complete, verified design together with an honest account of
what happens when that design is handed to an optimiser. Three results are worth
stating in advance:

\begin{itemize}\itemsep2pt
\item The controller is implemented twice, independently, and the two
implementations are shown to agree --- so the reported behaviour is a property
of the design rather than of one program.
\item Genetic tuning improves accuracy by 32.6\% and decisively beats a
budget-matched random search, but it silently destroys two structural properties
of the controller that were never written into its objective.
\item On the benchmark functions no optimiser dominates, and the cell with the
largest apparent difference between two algorithms proves statistically
insignificant once the individual runs are examined rather than their means.
\end{itemize}

\subsection{Why activity and preference are inputs}

Two of the four inputs are not the obvious sensors, and both follow from the
assistive-care setting.

\textbf{Activity} matters because thermoregulation is not uniform across
residents. Reduced mobility means reduced metabolic heat production, and
several conditions affecting mobility --- spinal cord injury in particular ---
also impair vasomotor and sweating responses below the level of the lesion. A
resting occupant therefore needs a \emph{warmer} room than a general-purpose
thermostat would choose, and the shift is large enough to warrant a first-class
input rather than a fixed offset.

\textbf{Preference} matters because comfort is subjective and, in an assistive
setting, contested: what a building manager considers efficient and what a
resident considers comfortable are different quantities. Exposing preference as
an input keeps the resident in the control loop rather than optimising them out
of it. It is a standing dial the resident sets, not a transient request.

%==============================================================================
\section{Controller Design}
\label{sec:t2-design}

\subsection{Linguistic variables}

Table~\ref{tab:t2-variables} gives the four inputs and two outputs. Universes
extend past the comfort band in both directions so that a fault condition --- a
heating failure in winter, solar gain in summer --- remains inside the modelled
range rather than being clipped. Section~\ref{sec:t2-cost} shows this to be the
first property genetic tuning destroys.

\begin{table*}[t]
\caption{The six linguistic variables. Universes extend past the comfort band
in both directions so that a fault condition --- a heating failure in winter,
solar gain in summer --- remains inside the modelled range rather than being
clipped at the boundary.}
\label{tab:t2-variables}
\centering
\small
\begin{tabular}{@{}lp{0.12\textwidth}p{0.20\textwidth}p{0.25\textwidth}p{0.23\textwidth}@{}}
\toprule
& \textbf{Variable} & \textbf{Universe} & \textbf{Fuzzy sets} &
\textbf{Sensed by} \\
\midrule
input  & \texttt{room\_temp}  & 14--34\,$^{\circ}$C
       & Cold, Cool, Comfortable, Warm, Hot
       & wall thermistor \\
input  & \texttt{activity}    & 0--10 index
       & Resting, Light, Active
       & PIR motion sensor + wrist accelerometer \\
input  & \texttt{daylight}    & 0--100\% of design lux
       & Dark, Dim, Bright
       & ceiling lux sensor \\
input  & \texttt{preference}  & $-5$ (cooler) \ldots\ $+5$ (warmer)
       & Cooler, Neutral, Warmer
       & resident-set dial, persisted \\
\midrule
output & \texttt{hvac}        & $-100$ \ldots\ $+100$\% of plant capacity
       & CoolHigh, CoolLow, Off, HeatLow, HeatHigh
       & signed actuator command \\
output & \texttt{dimmer}      & 0--100\% of full lamp output
       & Off, Low, Medium, High
       & lamp dimmer command \\
\bottomrule
\end{tabular}
\end{table*}

The HVAC output is a single \emph{signed} actuator rather than separate heating
and cooling channels. The two are physically exclusive, and a controller able to
command both at once can fight itself; a signed axis makes that state
unrepresentable rather than merely discouraged.

\subsection{Inference model: Mamdani rather than Sugeno}

Mamdani inference~\cite{mamdani1975} is used in preference to the
Takagi--Sugeno model~\cite{takagi1985}, for three reasons specific to this
application.

First, the rule base is elicited from human domain knowledge rather than fitted
to logged data. Mamdani's fuzzy consequents express that knowledge in the form
it arrives in --- ``if the room is cold and the resident is resting, heat
hard'' --- whereas a Sugeno consequent is a function of the inputs and has no
natural reading until it is evaluated.

Second, and decisively in an assistive-care setting, the behaviour must be
approved by someone who is not an engineer. An occupational therapist or carer
can check the sentence \texttt{IF room\_temp is Cold AND activity is Resting
THEN hvac is HeatHigh} against their own judgement. The first-order Sugeno
equivalent, $\texttt{hvac} = 71.2 - 3.4T + 1.1A + 0.02D + 4.6P$, is not
something a domain expert can audit, and after tuning its coefficients carry no
particular meaning at all. Interpretability is a functional requirement here,
not a nicety.

Third, the plant is slow: room air temperature moves over minutes, so Sugeno's
cheaper arithmetic buys nothing that matters at run time. Where it \emph{would}
matter is inside the tuning loop of Section~\ref{sec:t2-ga}, and
Section~\ref{sec:t2-sugeno} quantifies exactly what a Sugeno formulation would
change there.

\subsection{Membership functions}

Triangles are used for interior sets and trapezoids at the edges of every
universe (Fig.~\ref{fig:t2-mfs}).

Interior sets need a single unambiguous prototype value, which the apex of a
triangle provides: there is exactly one temperature that is most
\emph{Comfortable}. Edge sets must instead \emph{saturate}, because
8\,$^{\circ}$C and 14\,$^{\circ}$C are not meaningfully different to an occupant
and both should command full heat. A trapezoid with $a = b$ expresses that; a
triangle cannot, since it would make one extreme of the universe less Cold than
a point just inside it. Both shapes are piecewise linear, and so far cheaper to
evaluate than Gaussians or generalised bell curves --- immaterial for a single
inference, but not for the tuning loop, which evaluates the controller millions
of times.

Adjacent sets overlap at roughly half the base width, crossing near
$\mu = 0.5$~\cite{lee1990}. This buys two properties, both verified rather than
assumed in Section~\ref{sec:t2-impl}: the control surface stays continuous, and
no input value falls through the rule base.

\begin{figure}[t]
\centering
\includegraphics[width=\columnwidth]{fig01_membership_functions.png}
\caption{Membership functions for all six variables. Triangles inside,
trapezoids at the edges so the extremes saturate; adjacent sets cross near
$\mu = 0.5$.}
\label{fig:t2-mfs}
\end{figure}

\subsection{Rule base}

The controller has \textbf{54 rules}: 45 thermal rules over
$\texttt{room\_temp} \times \texttt{activity} \times \texttt{preference}$, and 9
lighting rules over $\texttt{daylight} \times \texttt{activity}$.

They are generated from additive fuzzy associative memory (FAM)
tables~\cite{kosko1992} rather than enumerated individually. Each antecedent
contributes an integer shift on a five-point demand scale; the shifts are summed
and the result saturated back into range:
\begin{align}
m &= \mathrm{clip}\big(s_{\text{act}} + s_{\text{pref}},\; -1,\; +1\big),
     \label{eq:modifier}\\
d &= \mathrm{clip}\big(b_{\text{temp}} + m,\; -2,\; +2\big),
\end{align}
with the demand $d$ mapped to \{CoolHigh, CoolLow, Off, HeatLow, HeatHigh\}.

The three lookup tables are small. Temperature contributes a base demand
$b_{\text{temp}} \in \{+2, +1, 0, -1, -2\}$ running from \texttt{Cold} to
\texttt{Hot}; activity contributes $s_{\text{act}} = +1$ for \texttt{Resting},
$0$ for \texttt{Light}, $-1$ for \texttt{Active}; preference contributes
$s_{\text{pref}} = -1, 0, +1$ for \texttt{Cooler}, \texttt{Neutral},
\texttt{Warmer}. To take one rule as a worked example, a \texttt{Cool} room
($b_{\text{temp}} = +1$) with an \texttt{Active} occupant
($s_{\text{act}} = -1$) at \texttt{Neutral} preference
($s_{\text{pref}} = 0$) gives $m = \mathrm{clip}(-1, -1, +1) = -1$ and
$d = \mathrm{clip}(1 - 1, -2, +2) = 0$, which maps to \texttt{Off}. That is
rule R17, and it is one of the four that fires in the worked scenario of
Section~\ref{sec:t2-behaviour}.

The additive structure makes the rule base \emph{auditable}. A carer can be
shown three small tables (Fig.~\ref{fig:t2-fam}) instead of a 45-row list, and
can check a claim such as ``resting shifts demand one step towards heating'' on
its own, rather than having to trust 45 independently asserted rules to be
mutually consistent. It also makes the rule base cheap to revise: changing the
metabolic assumption about resting occupants means editing one number, not
auditing fifteen rules for consistency afterwards.

A second property follows from the arithmetic rather than from inspection.
Because $b_{\text{temp}}$ is monotone in temperature and the modifier $m$ does
not depend on temperature at all, the demand $d$ is monotone in temperature for
every fixed activity and preference. The controller therefore cannot respond to
a colder room by heating less --- not merely does not, on the states that were
checked, but cannot, for any of the 45 combinations. Monotonicity of this kind
is normally an emergent property one has to test for; here it is a consequence
of how the rule base is constructed.

\textbf{The inner clip in \eqref{eq:modifier} carries a safety property.}
Without it, activity and preference each contribute a full step and can
therefore \emph{overrule} the temperature reading rather than adjust it: a
Comfortable 22.5\,$^{\circ}$C room with a resting occupant who preferred warmth
would sum to $+2$ and command \texttt{HeatHigh}, driving that room towards
30\,$^{\circ}$C. Clamping the combined modifier to a single step guarantees that
the measured temperature sets the demand to within one step, so the controller
can never heat a Hot room or cool a Cold one. That property holds by
construction, for every one of the 45 thermal rules, rather than by inspection.

\begin{figure*}[t]
\centering
\includegraphics[width=\textwidth]{fig15_fam_heatmaps.png}
\caption{The complete rule base as FAM matrices: three temperature $\times$
activity slices, one per preference setting, plus the lighting table. Colour is
the signed demand step, so the monotonicity the additive structure guarantees
appears as a gradient that never reverses along either axis --- the controller
cannot respond to a colder room by heating less. Saturation is visible at the
corners: Cold + Resting + Warmer and Cold + Resting + Neutral both command
\texttt{HeatHigh}, because once demand is at full heat there is nowhere further
to go.}
\label{fig:t2-fam}
\end{figure*}

The lighting rules follow the same principle on a four-point scale: activity
sets the illuminance the task needs, and available daylight is subtracted from
it, since the lamp need only make up the shortfall. \textbf{One entry departs
from the arithmetic.} A resting occupant in partial daylight (\texttt{Dim})
computes to needing no lamp at all. Falls are the dominant injury risk in this
population and a principal reason assistive-care housing exists, so the
controller maintains a floor of orientation lighting whenever daylight is not
\texttt{Bright}: $(\texttt{Dim}, \texttt{Resting}) \rightarrow \texttt{Low}$
rather than \texttt{Off}. This trades a small amount of energy for a safety
margin, and is the one place where the rule base is not a pure function of the
FAM arithmetic.

\subsection{Inference and defuzzification}

Conjunction is \texttt{min}, implication is \texttt{min} (clipping), aggregation
is \texttt{max}~\cite{mamdani1975,lee1990}. Each of the three is a choice, and
each was made for a stated reason rather than taken as a default.

\textbf{AND as \texttt{min} rather than \texttt{prod}.} The minimum t-norm makes
a rule's firing strength equal to its weakest satisfied antecedent, which is the
reading a domain expert intends by ``and'': a rule about a cold room with a
resting occupant should fire to the degree that the \emph{least} well satisfied
of those two conditions holds. The product t-norm instead multiplies, so three
antecedents each satisfied to degree 0.5 would fire at 0.125 rather than 0.5.
With three antecedents on every thermal rule, that would systematically
suppress the rules that fire in the overlap regions --- precisely the regions
where the controller is interpolating and where its output matters most.

\textbf{Implication as clipping rather than scaling.} A clipped consequent
retains a flat top at its rule's firing strength, which keeps each rule's
contribution individually visible in the aggregated set --- the plateaus in
Fig.~\ref{fig:t2-stages} can each be attributed to a named rule. Under scaling
the contributions blend into one another and the aggregate becomes difficult to
read, losing the property that motivated Mamdani inference in the first place.

\textbf{Aggregation as \texttt{max}.} Union rather than sum, so that two rules
recommending the same action do not produce a stronger recommendation than
either alone. With an additive FAM this matters: several rules routinely share a
consequent, and under summation their agreement would inflate the aggregated set
and bias the centroid, turning ``two rules agree'' into ``act twice as hard''.

The controller defuzzifies by centroid. All three standard candidates were
implemented and compared on measured behaviour rather than convention
(Fig.~\ref{fig:t2-defuzz}). Centroid uses every rule's contribution, which is
the point of having run the inference, and moves smoothly as the inputs move.
The bisector is cheaper and less sensitive to a long thin tail, but advances in
steps as the equal-area crossing jumps between grid cells --- behaviour an
actuator should not be asked to follow. Mean-of-maximum reaches further towards
the ends of the universe but ignores every set that is not the strongest,
discarding most of the information the rule base has just produced, and jumps as
the dominant rule changes~\cite{leekwijck1999}.

The comparison also quantifies a real cost of that choice, discussed in
Section~\ref{sec:t2-behaviour}: centroid cannot command the full actuator range.

\begin{figure*}[t]
\centering
\includegraphics[width=\textwidth]{fig16_defuzzifier_comparison.png}
\caption{Defuzzification compared on measured behaviour. \textbf{Left:} the
three methods applied to the same aggregated set. \textbf{Right:} the range of
HVAC commands each can actually reach, measured over a 70{,}875-point input
grid. Centroid spans only 79.9\% of the $\pm 100\%$ universe against
mean-of-maximum's 90.0\% --- a real limit on actuator authority, and the price
paid for smoothness.}
\label{fig:t2-defuzz}
\end{figure*}

%==============================================================================
\section{Implementation and Verification}
\label{sec:t2-impl}

\subsection{Two independent implementations}

The controller exists in two forms: a purpose-written Python Mamdani inference
engine and a custom base-MATLAB implementation. The MATLAB controller definition
is generated from the same variables, rules and membership functions. This
reduces transcription errors; agreement must still be checked numerically.

This yields a genuine verification opportunity. The two are identical in
\emph{definition} but independent in \emph{arithmetic}: different
implementations of fuzzification, clipping, aggregation and centroid
integration. Agreement between them is therefore evidence that both are
correct, and disagreement would localise a defect to one of them. The inference
engine is separately validated against an established open-source fuzzy
library, which must reproduce every consequent activation to within $10^{-12}$
and must converge on the same crisp output as the output universe is refined.

Two design decisions in the MATLAB implementation are worth stating because
they are the ones that make the comparison meaningful rather than
approximate.

\textbf{The MATLAB system is generated, not transcribed.} Maintaining two
hand-written definitions of one controller and hoping they stay in step is
exactly the kind of silent drift that invalidates a cross-implementation check:
the two could eventually differ, and the difference could be mistaken for a
finding. Regeneration and cross-implementation checks reduce this risk.

\textbf{Mamdani inference is implemented directly rather than called from a
library.} The controller is expressed as a plain data structure and the five
inference stages are evaluated explicitly. This was originally a constraint ---
the MATLAB licence available for this work does not include the fuzzy inference
library, so those functions were unavailable. The custom engine exposes firing
strengths, clipped consequents and aggregated sets to support the worked
examples. Plots are derived from these intermediate values and checked against
the numerical trace. This choice does not constitute use of the named Toolbox.

\subsection{Base-MATLAB implementation}
\label{sec:t2-matlab}

Figures~\ref{fig:t2-matlab1}, \ref{fig:t2-matlab2} and~\ref{fig:t2-matlab3} show
the MATLAB visualisations: membership functions, rule inference for a worked
scenario, and control surfaces. They are not evidence of a Fuzzy Logic Toolbox
implementation. The saved code uses a custom engine because that toolbox was
unavailable on the project's licence. The brief permits a simulated system
programmed in MATLAB, FuzzyLite or Juzzy and provides Toolbox reference links.
The base-MATLAB implementation follows that programming-language option without
claiming to use the Toolbox.

The evaluation reports a direct numerical comparison against the reference
engine. Both custom implementations use 501 output samples: unequal grids could
change centroid estimates independently of inference correctness. The recorded
base-MATLAB code check was executed in GNU Octave, with approximately
$3.6\times10^{-7}$ difference on the worked scenario, within the $10^{-3}$
tolerance on both actuators. This is separate from the Python/scikit-fuzzy tests.

\begin{figure}[t]
\centering
\includegraphics[width=\columnwidth]{matlab_fig1_membership.png}
\caption{MATLAB visualisation: membership functions for all four inputs
and both outputs. Compare with Fig.~\ref{fig:t2-mfs}, produced independently by
the reference engine.}
\label{fig:t2-matlab1}
\end{figure}

\begin{figure}[t]
\centering
\includegraphics[width=\columnwidth]{matlab_fig2_ruleinference.png}
\caption{MATLAB rule inference for the worked scenario
($\texttt{room\_temp} = 21$, $\texttt{activity} = 7.5$,
$\texttt{daylight} = 25$, $\texttt{preference} = 0$): the firing strengths, and
the aggregated output set with its centroid for each actuator.}
\label{fig:t2-matlab2}
\end{figure}

\begin{figure}[t]
\centering
\includegraphics[width=\columnwidth]{matlab_fig3_surfaces.png}
\caption{MATLAB control surfaces for both actuators. Compare with
Fig.~\ref{fig:t2-surface1}.}
\label{fig:t2-matlab3}
\end{figure}

\subsection{Structural audit}

The controller was swept over a 70{,}875-point grid covering the whole input
space (Table~\ref{tab:t2-audit}).

\begin{table}[t]
\caption{Structural audit over a 70{,}875-point input grid. Each row is a
property the design intends, measured rather than assumed.}
\label{tab:t2-audit}
\centering
\small
\begin{tabular}{@{}lr@{}}
\toprule
\textbf{Property} & \textbf{Measured} \\
\midrule
Inputs with no defined output    & \textbf{0} \\
Dead rules (never fire)          & \textbf{0} \\
Rules active per input, mean     & 2.89 \\
Rules active per input, max      & 12 \\
Inferences per second            & 15{,}729 \\
Centroid, exact vs.\ discrete sum & $3.6 \times 10^{-15}$ \\
\bottomrule
\end{tabular}
\end{table}

Zero uncovered inputs and zero dead rules are the two properties the overlap
choice was made to buy, and both hold. Figure~\ref{fig:t2-coverage} shows the
distributions behind those two counts. The left panel is the number of rules
firing simultaneously: it has no zero bin, so no input state falls through the
rule base. The right panel is how often each of the 54 rules fires across the
space; no bar is zero, so no rule is dead. Both are consequences of overlapping
adjacent sets near $\mu = 0.5$, and both are the kind of claim that is cheap to
assert and easy to get wrong, so they are measured rather than assumed.

The distribution is also informative about the cost of inference. A mean of
2.89 simultaneously active rules out of 54 means the controller is doing very
little work per decision: a reading typically lands in two sets on temperature
and one or two on the other variables, and the rule base is sparse enough that
only a handful of combinations survive. The maximum of 12 occurs where a
reading sits on a crossover point of three antecedent variables at once. That
sparsity is what makes 15{,}729 inferences per second achievable without any
optimisation, and it is a direct consequence of the overlap choice rather than
of implementation effort.

The final row of Table~\ref{tab:t2-audit} confirms that the integration scheme
introduces no error of its own: exact trapezoidal integration of the aggregated
set and a discrete centre-of-gravity sum over the same grid agree to
$3.6 \times 10^{-15}$ on the HVAC output, which is machine precision. Any
disagreement between this controller and another implementation is therefore
attributable to the model, not to the quadrature.

\begin{figure}[t]
\centering
\includegraphics[width=\columnwidth]{fig19_rule_coverage.png}
\caption{Rule-base coverage over 20{,}000 uniformly sampled input states.
\textbf{Left:} how many rules fire simultaneously --- the distribution has no
zero bin, so no input state falls through the rule base. \textbf{Right:} how
often each of the 54 rules fires, thermal rules and lighting rules
distinguished; no bar is zero, so there are no dead rules.}
\label{fig:t2-coverage}
\end{figure}

%==============================================================================
\section{Analysis of Output Behaviour}
\label{sec:t2-behaviour}

\subsection{A worked inference, stage by stage}

Consider an evening scenario: the resident is moving about preparing a meal as
the light fades. $\texttt{room\_temp} = 21\,^{\circ}$C,
$\texttt{activity} = 7.5$, $\texttt{daylight} = 25\%$, $\texttt{preference} = 0$.

\textbf{Fuzzification} (Fig.~\ref{fig:t2-fuzzification}) gives
$\mu_{\text{Cool}} = 0.40$ and $\mu_{\text{Comfortable}} = 0.25$ on temperature,
$\mu_{\text{Active}} = 0.50$ on activity, $\mu_{\text{Dark}} = 0.25$ and
$\mu_{\text{Dim}} = 0.20$ on daylight, and $\mu_{\text{Neutral}} = 1.0$ on
preference.

\textbf{Rule activation} (Fig.~\ref{fig:t2-activation}) fires four of the 54
rules. R16 (\texttt{Cool} $\wedge$ \texttt{Active} $\wedge$ \texttt{Neutral}
$\rightarrow$ \texttt{Off}) at $\alpha = 0.40$ and R25 (\texttt{Comfortable}
$\wedge$ \texttt{Active} $\wedge$ \texttt{Neutral} $\rightarrow$
\texttt{CoolLow}) at $\alpha = 0.25$ carry the thermal decision; two lighting
rules carry the lamp.

\textbf{Implication and aggregation} (Fig.~\ref{fig:t2-stages}) clip each
consequent set at its rule's firing strength and take the union.

\textbf{Defuzzification} returns $\texttt{hvac} = -22.77\%$ --- gentle cooling,
since the occupant is generating heat in an already-adequate room --- and
$\texttt{dimmer} = 72.60\%$, substantial light for a task in falling daylight.
Both are the actions a carer would expect, which is the qualitative check that
matters most.

\begin{figure}[t]
\centering
\includegraphics[width=\columnwidth]{fig02_fuzzification_worked.png}
\caption{Fuzzification of the worked scenario: each crisp reading dropped onto
its variable's sets, with the resulting degrees of membership marked.}
\label{fig:t2-fuzzification}
\end{figure}

\begin{figure}[t]
\centering
\includegraphics[width=\columnwidth]{fig03_rule_activation.png}
\caption{Rule activation for the worked scenario: 4 of 54 rules fire, with their
firing strengths. Sparse activation is the intended consequence of overlapping
adjacent sets at $\mu \approx 0.5$ rather than more widely.}
\label{fig:t2-activation}
\end{figure}

\begin{figure}[t]
\centering
\includegraphics[width=\columnwidth]{fig04_inference_stages_hvac.png}
\caption{Implication, aggregation and defuzzification for \texttt{hvac}. Each
consequent set is clipped at its rule's $\alpha$ (the flat tops), the union is
taken, and the centroid of the result is the commanded value. The same three
stages for the \texttt{dimmer} output are in Fig.~\ref{fig:t2-stages-dimmer} in
the appendix.}
\label{fig:t2-stages}
\end{figure}

\subsection{Control surfaces}

A control surface is the whole input--output map of the controller in one
picture, and it is where design intent either shows up or does not.
Figure~\ref{fig:t2-surface1} gives the HVAC actuator against temperature and
activity; Fig.~\ref{fig:t2-surface2} gives the dimmer against daylight and
activity. Figure~\ref{fig:t2-surface-pref} in the appendix gives the third
slice, HVAC against temperature and preference.

Four intended properties are visible, and each corresponds to a design decision
made in Section~\ref{sec:t2-design}.

\textbf{Temperature dominates.} The HVAC surface falls monotonically from full
heat to full cool as the room warms, and no combination of activity and
preference reverses that. This is the clamped modifier of \eqref{eq:modifier}:
because the combined activity-plus-preference shift is capped at one step on a
five-step scale, the measured temperature always sets the demand to within one
step. Monotonicity here is a structural guarantee of the rule base rather than
an emergent property that happened to hold.

\textbf{Activity and preference shift the surface without tilting it.} Moving
along the activity axis translates the whole temperature response rather than
changing its slope, which is exactly the intended semantics: a resting occupant
should be given a warmer room, not a differently-shaped response.
Figure~\ref{fig:t2-response} in the appendix shows the same effect as three
overlaid cross-sections, where the ordering is easier to read off than from a
three-dimensional surface.

\textbf{The surface is continuous, with no cliffs.} Adjacent sets overlap near
$\mu = 0.5$, so as a reading crosses from one set into the next the two
consequents blend rather than switch. An actuator can follow this; it would not
be able to follow a staircase. The visible terracing is the five-level demand
scale, softened by the overlap into ramps rather than steps.

\textbf{The dimmer surface makes up a shortfall.} The lamp output falls as
daylight rises and rises with activity, which is the subtraction the lighting
FAM performs. The flat floor at low output is the fall-risk override: the lamp
never switches off entirely while daylight is below \texttt{Bright}, which is
visible as a plateau that does not descend to zero.

\begin{figure}[t]
\centering
\includegraphics[width=\columnwidth]{fig06_surface_hvac_temp_activity.png}
\caption{Control surface: \texttt{hvac} against \texttt{room\_temp} and
\texttt{activity}, at neutral preference. Monotone in temperature, shifted but
not tilted by activity, and continuous throughout. The saturation plateaus at
roughly $\pm 79.9\%$ are the centroid authority limit discussed below.}
\label{fig:t2-surface1}
\end{figure}

\begin{figure}[t]
\centering
\includegraphics[width=\columnwidth]{fig08_surface_dimmer.png}
\caption{Control surface: \texttt{dimmer} against \texttt{daylight} and
\texttt{activity}. The lamp makes up the shortfall between what the task needs
and what daylight supplies. The floor at low output is the fall-risk override:
orientation lighting is never fully extinguished while daylight is below
\texttt{Bright}.}
\label{fig:t2-surface2}
\end{figure}

\subsection{An operational day}

A surface shows what the controller \emph{can} do; a time series shows what it
actually does in the states a flat occupies. Figure~\ref{fig:t2-day} runs a
24-hour profile through the controller --- a cold night, a resident waking and
transferring, a mild sunlit afternoon, an active evening as daylight fades.

Over the simulated day the HVAC command spans $[-70.97, +76.60]\%$ and the
dimmer $[7.0, 87.69]\%$, with 12.8 heating hours and 6.5 cooling hours. Three
behaviours the design intends are visible in the traces. The lamp never reaches
zero, which is the fall-risk floor holding across every state the day visits.
The HVAC command never changes sign abruptly: it passes through the
\texttt{Off} region rather than switching between heating and cooling, which is
the behaviour a real plant needs if it is not to short-cycle. And the morning
transfer --- a brief burst of high activity in a cool room --- reduces the heat
demand rather than raising it, because the occupant is generating metabolic
heat, which is the assistive-care-specific behaviour that motivated making
activity an input at all.

\begin{figure}[t]
\centering
\includegraphics[width=\columnwidth]{fig10_operational_day.png}
\caption{A simulated operational day. Sensor inputs above, commanded actions
below. The dimmer never reaches zero and the HVAC command transitions smoothly
rather than switching.}
\label{fig:t2-day}
\end{figure}

\subsection{Two limitations the analysis exposes}

Both of the following emerged from analysing output behaviour rather than from
inspecting the design, which is an argument for performing this analysis at all.

\textbf{The controller cannot reach full actuator authority.} Centroid
defuzzification spans only $[-79.9, +79.9]$ of the $\pm 100\%$ HVAC universe,
because the aggregated set always retains area on the interior side and the
centre of gravity is pulled inwards. Measured against bisector (81.2\%) and
mean-of-maximum (90.0\%) in Fig.~\ref{fig:t2-defuzz}, this is a real limit
rather than a plotting artefact: commanding 100\% capacity is impossible
whatever the inputs. It is acceptable here because the plant is oversized
relative to demand and because the alternatives trade it for discontinuity, but
a deployment requiring full authority would need to rescale the output universe.

\textbf{The HVAC response is not pointwise monotone in temperature}, by about
0.6\% of the output span. Where the demand scale saturates, two adjacent
temperature sets share a consequent, and that consequent's activation is
V-shaped through their crossover. The effect is bounded and an order of
magnitude below any real actuator's resolution.

%==============================================================================
\section{Genetic Tuning of the Membership Functions}
\label{sec:t2-ga}

The controller structure of Section~\ref{sec:t2-design} is now held fixed ---
the same four inputs, two outputs and 54 rules --- and only the membership
functions are tuned against a set of input--output examples, following the
genetic-fuzzy approach surveyed by Cord\'{o}n et al.~\cite{cordon2004}.

\subsection{The dataset}

\textbf{The training data is synthetic, and nothing derived from it is evidence
about real occupants.} No public dataset of assistive-care environmental
preferences exists, and collecting one would require ethical approval and
resident participation. The examples are generated instead from an explicitly
stated analytical reference policy --- saturating linear ramps on setpoint error
and on task illuminance --- with Gaussian noise at 3\% of each output's span
representing annotator disagreement. 400 examples are drawn and split 280
training / 120 held-out test.

Two properties of that construction matter. Inputs are sampled \emph{uniformly}
rather than from a plausible occupancy distribution: a controller tuned only on
the states a flat usually occupies would leave the membership functions covering
rare states --- a heating failure, a heatwave --- entirely untuned, and those are
the states where the controller matters most. And the reference policy is
deliberately a \emph{different functional form} from the controller. Were the
target generated by a Mamdani system with the same rule base, the error could in
principle be driven to zero, demonstrating only that the encoding can represent
its own output. Fitting a genuinely different shape is the harder test, and
residual error at convergence is therefore expected rather than a failure.

Accordingly, the improvement reported below establishes that the search can
recover a known target. It does not establish that the tuned controller would
serve real residents better.

\subsection{Problem encoding and chromosome length}

Encoding is \textbf{binary}: one gene of 8 bits per tunable membership-function
parameter, decoded linearly onto the physical range of the variable it belongs
to,
\begin{equation}
v = \ell + \frac{b}{255}\,(h - \ell),
\end{equation}
so a room-temperature breakpoint decodes somewhere in $[14, 34]$ degrees and an
HVAC breakpoint somewhere in $[-100, +100]$ percent. A fixed word length makes
the chromosome length a stated, checkable quantity~\cite{goldberg1989}.

The parameter layout, and hence the \textbf{chromosome length}, is:

\begin{center}
\small
\begin{tabular}{@{}lr@{\quad}lr@{}}
\texttt{room\_temp}  & 17 & \texttt{hvac}   & 17 \\
\texttt{activity}    & 11 & \texttt{dimmer} & 14 \\
\texttt{daylight}    & 11 & & \\
\texttt{preference}  & 11 & & \\
\multicolumn{4}{c}{\textbf{81 parameters $\times$ 8 bits = 648 bits}} \\
\end{tabular}
\end{center}

The resolution this buys is $(h-\ell)/255$: about $0.078\,^{\circ}$C on room
temperature and 0.78 percentage points on the HVAC command, finer than either
sensor or actuator resolves. Quantisation is therefore not what limits the
result, though it is not free: encoding and decoding the hand-designed
controller displaces its parameters by at most 0.392 units and costs
$+0.000342$ in normalised error before the search gains anything --- roughly
0.7\% of what the search subsequently recovers.
Figure~\ref{fig:t2-encoding} in the appendix breaks that cost down per gene.

\textbf{Infeasible chromosomes are repaired, not rejected.} A decoded chromosome
routinely proposes a triangle whose parameters are out of order. Discarding the
individual would waste the evaluation and bias the search towards whatever
region happens to be feasible, so each membership function's parameters are
sorted ascending on decode. Every point in the 648-bit space therefore maps to a
valid controller, and selection pressure stays on fitness.

\subsection{Genetic operators}

Roulette-wheel (fitness-proportionate) selection with elitism 2; one-point
crossover at $p_c = 0.7$; per-bit mutation at $p_m = 0.01$; population 60 over
150 generations, giving 9{,}060 evaluations per run~\cite{goldberg1989}.

The crossover cut may fall anywhere in the chromosome, \emph{including inside a
gene}. Cutting only on gene boundaries would make the 8-bit words atomic and
prevent the search from refining a single breakpoint by a small amount. Elitism
guarantees the best individual survives untouched, so best fitness cannot fall
between generations.

The population is \textbf{seeded with the hand-designed controller}, which makes
``never worse than the expert design'' a property of the search rather than a
hope. The seeded individual is charged an evaluation like any other.

\subsection{Fitness function}

Span-normalised RMSE against the reference data, expressed as a quantity to
maximise:
\begin{equation}
\mathrm{fitness} = \frac{1}{1 + \mathrm{nRMSE}},
\qquad
\mathrm{nRMSE} = \frac{1}{2}\sum_{j=1}^{2}
\frac{\sqrt{\overline{(\hat{y}_j - y_j)^2}}}{h_j - \ell_j}.
\end{equation}

Two choices inside that are worth stating. \textbf{Normalising by span} is
necessary because the outputs live on different scales: HVAC spans 200 units of
signed capacity, the dimmer 100 units of percent. Averaging raw squared error
would let HVAC dominate roughly four to one, leaving the lighting membership
functions effectively untuned. \textbf{Undefined outputs are charged the full
span} rather than dropped, so the search cannot improve its score by evolving a
controller with holes in it --- a provision that Section~\ref{sec:t2-cost} shows
discourages holes without preventing them.

Roulette-wheel selection requires a positive quantity that is larger for better
individuals; $1/(1+e)$ is bounded in $(0, 1]$ and monotone in the error, so it
changes which individuals are selected without changing which is best.

\subsection{Results against a budget-matched control}

A genetic algorithm that improves on its starting point has not thereby shown
that \emph{evolution} did the work: a search that merely samples the space
finely can do the same. Separating the two requires spending an identical number
of fitness evaluations sampling at random. Both searches are seeded with the
expert controller, both are charged 9{,}060 evaluations per run, both use the
same fitness function, and both are repeated over five seeds.

\begin{table*}[t]
\caption{Held-out test error over five seeds. Both searches were given an
identical budget of 9{,}060 fitness evaluations, an identical expert-seeded
starting point and the identical fitness function; they differ in exactly one
respect, which is how candidates are proposed. Lower error is better.}
\label{tab:t2-ga}
\centering
\small
\begin{tabular}{@{}lrrrrrl@{}}
\toprule
\textbf{Method} & \textbf{mean} & \textbf{std} & \textbf{best} &
\textbf{worst} & \textbf{vs.\ baseline} & \textbf{note} \\
\midrule
Hand-designed baseline   & 0.13299 & ---     & ---     & ---     & ---
& the expert design, untuned \\
Random search (matched)  & 0.13340 & 0.00000 & 0.13340 & 0.13340 & $-0.3\%$
& never improved on its seed \\
\textbf{Genetic algorithm} & \textbf{0.08964} & 0.00438 & \textbf{0.08504}
& 0.09682 & \textbf{$+32.6\%$} & every seed beat the baseline \\
\bottomrule
\end{tabular}
\end{table*}

The genetic algorithm improves on the hand design by \textbf{32.6\%} and beats
the control decisively (Table~\ref{tab:t2-ga}, Fig.~\ref{fig:t2-control}).

The convergence traces (Fig.~\ref{fig:t2-gaconv}) add two details the summary
table hides. First, every one of the five seeds improves on the baseline, and
they do so at markedly different rates --- one seed reaches its final quality by
generation 40 while another is still descending at generation 140. That spread
is ordinary for a stochastic search, and it is the reason the comparison is run
over five seeds rather than one. Second, and more usefully, \emph{none} of the
five has clearly plateaued by the 150-generation budget. The search is still
finding improvements when it is stopped, which says the binding constraint is
the evaluation budget rather than the encoding or the operators. A longer run
would be expected to do better still; whether that would be a good idea is a
separate question, and Section~\ref{sec:t2-cost} argues it would not.

Figure~\ref{fig:t2-mfba} shows what the tuning did to the membership functions
themselves. The sets migrate substantially --- this is not fine polishing of the
expert design --- and on several variables they visibly stop meeting the ends of
their universes, which is the defect Section~\ref{sec:t2-cost} quantifies.
Figure~\ref{fig:t2-predtarget} in the appendix shows the same improvement in the
output space, as predictions against targets on the held-out split.

\begin{figure}[t]
\centering
\includegraphics[width=\columnwidth]{fig11_ga_convergence.png}
\caption{Convergence of the genetic algorithm over five seeds: best training
error against generation, with the hand-designed baseline marked. Every seed
improves on the baseline; none has clearly plateaued at the 150-generation
budget, which indicates the budget rather than the encoding is the binding
constraint.}
\label{fig:t2-gaconv}
\end{figure}

\textbf{Random search did not improve on its seed even once.} Its result is
identical to five decimal places across all five runs, with standard deviation
exactly zero, and that value is precisely the expert chromosome it was seeded
with. In $5 \times 9{,}060 = 45{,}300$ uniform draws from a 648-bit space, not
one beat the hand design. The mechanism is plain: the space has $2^{648}$
points, and uniform sampling of it is hopeless in a way that uniform sampling of
a three-parameter space is not. Dimensionality, not sophistication, is what
determines whether the evolutionary machinery earns its keep here.

\begin{figure*}[t]
\centering
\includegraphics[width=\textwidth]{fig17_part2_control.png}
\caption{The budget-matched control. \textbf{Left:} held-out test error. The
five random-search points coincide exactly, because all five runs returned the
expert chromosome they were seeded with. \textbf{Right:} convergence against the
level random search reached --- a flat line, since it never left its seed.}
\label{fig:t2-control}
\end{figure*}

\begin{figure*}[t]
\centering
\includegraphics[width=.82\textwidth,height=.72\textheight,keepaspectratio]{fig12_mf_before_after.png}
\caption{Membership functions before (left column) and after (right column)
genetic tuning, for the best of the five runs. All six variables are shown. The
sets migrate substantially rather than being finely polished, and on several
variables the outermost sets no longer reach the ends of their universe --- the
coverage failure quantified in Section~\ref{sec:t2-cost}. Compare the left
column with Fig.~\ref{fig:t2-mfs}: the hand design places a saturating trapezoid
at each end precisely so that the extremes are fully covered, and that is the
property tuning removes first.}
\label{fig:t2-mfba}
\end{figure*}

\subsection{What the tuning cost}
\label{sec:t2-cost}

The 32.6\% improvement is real, and reporting it alone would be misleading. Two
further audits were run on the tuned controllers, and both fail.

\textbf{Linguistic ordering breaks in three of five runs.} Mamdani inference was
adopted because the rule base stays readable, and that argument survives
optimisation only if \texttt{Cold} remains to the left of \texttt{Cool}
afterwards. A controller whose sets have reordered has rules that no longer mean
what they say, and the interpretability justification collapses with it.
Comparing set centres of mass before and after (Fig.~\ref{fig:t2-ordering}), two
runs reorder \texttt{room\_temp} and one reorders \texttt{hvac}.

\textbf{Tuning punches holes in the rule base.} The hand-designed controller
returns a defined output for every one of 50{,}000 uniformly sampled input
states. Four of the five tuned controllers do not: between 0.37\% and 0.80\% of
states activate \emph{no rule at all}, leaving the aggregated set empty, the
centroid undefined, and the actuator without a command
(Fig.~\ref{fig:t2-holes}).

The mechanism is precise, and it bears directly on the design rationale of
Section~\ref{sec:t2-design}. Trapezoids were placed at the universe edges
specifically so the extremes would saturate. After tuning, the outermost sets no
longer reach those edges: in one run, \emph{every} input has $\mu = 0$ at both
ends of its universe. The states that fall through are therefore exactly the
fault conditions --- heating failure, heatwave, total darkness, maximum stated
preference --- that the universes were widened to capture. Unconstrained
optimisation removed the property the edge trapezoids existed to provide.

\begin{figure*}[t]
\centering
\includegraphics[width=\textwidth]{fig22_linguistic_ordering.png}
\caption{Does the tuned controller still mean what its rules say? Set centres of
mass before (left tick) and after (right tick) tuning, per run. A crossing is a
reordering: two runs reorder \texttt{room\_temp} and one reorders
\texttt{hvac}.}
\label{fig:t2-ordering}
\end{figure*}

\begin{figure*}[t]
\centering
\includegraphics[width=\textwidth]{fig20_error_maps.png}
\caption{\textbf{Left, centre:} where the controller disagrees with the
reference policy, before and after tuning. Tuning shrinks the error band rather
than removing it, as expected when the target is a different functional form.
Dark cells are states where no rule fires. \textbf{Right:} the coverage audit
over the full four-dimensional input space. The hand design has no undefined
state; four of five tuned controllers do.}
\label{fig:t2-holes}
\end{figure*}

Taken together (Table~\ref{tab:t2-audit2}), \textbf{no run produced a controller
that was both fully covering and linguistically ordered}. The two runs that kept
their ordering both lost coverage; the one run that kept full coverage reordered
its output sets.

\begin{table*}[t]
\caption{Per-run audit of the tuned controllers. Accuracy improved in every one
of the five runs; not one preserved both structural properties that the
hand-designed controller has. ``Coverage'' is the percentage of 50{,}000
uniformly sampled input states that activate no rule at all and therefore leave
the actuator without a command; ``ordering'' is whether every variable's fuzzy
sets are still in their original left-to-right order.}
\label{tab:t2-audit2}
\centering
\small
\begin{tabular}{@{}lrccl@{}}
\toprule
& \textbf{test error} & \textbf{states with no output} &
\textbf{ordering intact} & \textbf{verdict} \\
\midrule
Hand design & 0.13299 & 0.00\% & yes & both properties hold \\
\midrule
Run 1 & 0.09682          & 0.39\% & no  & fails both \\
Run 2 & \textbf{0.08504} & 0.80\% & yes & loses coverage \\
Run 3 & 0.08843          & 0.00\% & no  & loses ordering \\
Run 4 & 0.08980          & 0.37\% & no  & fails both \\
Run 5 & 0.08810          & 0.42\% & yes & loses coverage \\
\midrule
\multicolumn{4}{@{}l}{\textbf{Runs preserving both properties}} &
\textbf{0 of 5} \\
\bottomrule
\end{tabular}
\end{table*}

The lesson generalises beyond this controller. The fitness function encodes
accuracy and nothing else, so accuracy is what it delivered --- at the cost of
two safety-relevant structural properties that nobody decided to trade, because
nobody had written them into the objective. Both are cheap to enforce: ordering
by extending the repair operator to sort set centres as well as parameters,
coverage by a penalty term or by pinning the outermost breakpoints to the
universe edges. Neither was applied here, so that the question of whether
unconstrained tuning preserves interpretability could be answered by measurement
rather than assumed.

\subsection{The solution under a Sugeno model}
\label{sec:t2-sugeno}

Under Takagi--Sugeno inference only the \emph{consequent} side changes. The
antecedents, the input membership functions and all 54 rules are unaffected, so
the 50 input parameters of the chromosome are untouched. What changes is the 31
output parameters, and the range of outcomes is wide
(Table~\ref{tab:t2-sugeno}).

\begin{table}[t]
\caption{The chromosome under Mamdani and three Sugeno formulations.}
\label{tab:t2-sugeno}
\centering
\small
\begin{tabular}{@{}lrrl@{}}
\toprule
\textbf{Formulation} & \textbf{Params} & \textbf{Bits} &
\textbf{Auditable?} \\
\midrule
Mamdani (as built)          & 81  & \textbf{648}   & yes \\
Zero-order Sugeno           & 59  & 472            & partly \\
First-order Sugeno          & 320 & 2{,}560        & no \\
First-order + least squares & 50  & \textbf{400}   & no \\
\bottomrule
\end{tabular}
\end{table}

Under \textbf{zero-order} Sugeno each output fuzzy set becomes a single
singleton, so the 5 HVAC and 4 dimmer sets become 9 constants and the chromosome
shortens by 27\%. Under \textbf{first-order} Sugeno each of the 54 rules carries
its own linear function of four inputs, so the chromosome grows almost fourfold
to 2{,}560 bits and the search space from $2^{648}$ to $2^{2560}$. Given that
9{,}060 evaluations against 648 bits does not exhaust the present search, a
four-times-longer chromosome increases the search burden at the same budget;
its actual performance would require a separate experiment.

\textbf{The decisive gain is that least squares becomes available.} With the
antecedents fixed, a Sugeno system is \emph{linear in its consequent
parameters}, so the optimal coefficients can be solved directly rather than
searched. This enables the hybrid ANFIS-style scheme~\cite{jang1993}, which
would be the approach of choice: the genetic algorithm tunes only the 50 input
membership-function parameters (400 bits), and least squares solves the
consequents exactly at every fitness evaluation. The search space shrinks by
more than a third \emph{and} every individual is scored using optimal
consequents. This is precisely why ANFIS is built on Sugeno rather than Mamdani:
centroid defuzzification is not linear in the output membership-function
parameters, so no equivalent closed-form step exists for the system built here.

Evaluation would also become cheaper --- a weighted average over active rules
replaces integrating a 501-point aggregated set per output per sample --- and
the repair operator would apply only to the input half, since consequent
coefficients carry no ordering constraint.

What deteriorates is the property for which Mamdani was chosen.
The present output-set interpretation and linguistic-ordering audit of
Section~\ref{sec:t2-cost} would no longer apply directly. Sugeno rules can still
be inspected, but would require a different account of their output behaviour.
\textbf{More seriously, the safety property would need re-establishing.} ``Never
heat a hot room'' currently follows from the \emph{structure} of the clamped FAM.
With free linear consequents an optimiser could produce coefficients that
violate it while still scoring well on average error, so the property would have
to be reimposed as a constraint or a fitness penalty --- and a penalty only
discourages violations rather than preventing them, which is strictly weaker
than what the Mamdani system provides. Section~\ref{sec:t2-cost} is a
demonstration of how much weaker.

Sugeno offers a computationally attractive alternative. For this assistive-care
scenario, it would need explicit behavioural constraints and a suitable audit
before an accuracy gain could justify replacing the inspectable Mamdani design.

%==============================================================================
\section{Optimiser Comparison on CEC'2005 Functions}
\label{sec:t2-cec}

\subsection{Functions and reproducible instances}
Two functions from the CEC'2005 real-parameter suite~\cite{suganthan2005} provide contrasting
search landscapes. F6 is shifted Rosenbrock, a non-separable problem with a curved valley;
F9 is shifted Rastrigin, a separable problem with many local minima. Their definitions are
\begin{align}
F_6(\mathbf{x}) &= \sum_{i=1}^{D-1}[100(z_i^2-z_{i+1})^2+(z_i-1)^2]+390,\\
F_9(\mathbf{x}) &= \sum_{i=1}^{D}[z_i^2-10\cos(2\pi z_i)+10]-330.
\end{align}
For F6, $\mathbf{z}=\mathbf{x}-\mathbf{o}+1$, with domain $[-100,100]^D$ and optimum 390.
For F9, $\mathbf{z}=\mathbf{x}-\mathbf{o}$, with domain $[-5,5]^D$ and optimum $-330$.
Both use the first $D$ entries of the \emph{official} CEC2005 shift files, matching the reference
MATLAB implementation. The vendored data include source attribution and checksums; earlier
experiments using locally generated shifts are archived and are not the results reported here.
Figure~\ref{fig:t2-landscapes} visualises the two official two-dimensional landscapes.

There are 15 independent repetitions for each algorithm, function and dimension, with
$D\in\{2,10\}$: 180 runs in total. Every run receives $10,000D$ objective evaluations.
The random stream is determined by the repetition, algorithm, dimension and function; streams
are reproducible but do not represent identical initial populations across different algorithms.
An objective wrapper counts evaluations and rejects calls after the budget is exhausted.
The two functions complement each other, but success on them does not establish general
optimiser superiority across the CEC suite.

\subsection{Algorithms and evaluation}
The \textbf{GA} is real-coded, with population 50, tournament size three, blend crossover
($\alpha=0.5$) at probability 0.9, Gaussian mutation with scale 0.1 times the search range and
probability 0.1, and two-member elitism. Real coding avoids discretising these continuous
benchmarks; it differs from the deliberately binary chromosome in the controller experiment.

\textbf{Particle swarm optimisation (PSO)}~\cite{kennedy1995} uses a swarm of 40 particles,
global-best topology, inertia decreasing linearly from 0.9 to 0.4, and cognitive/social
coefficients 1.49445. Velocities are capped at 0.2 times the domain range and positions are
bounded. Its observed poor runs should not be read as a limit on other PSO variants or
attributed to absent velocity control: this implementation already has that control.

\textbf{Simulated annealing (SA)}~\cite{kirkpatrick1983} uses a single candidate with stochastic
proposals and temperature-dependent acceptance of worse moves. Temperature decreases by a
factor of 0.995 after each chain of 20 proposals; the Gaussian proposal scale decreases from
0.20 to 0.002 times the range. Initial temperature is calibrated to accept a typical uphill
move with probability 0.5, and there are no restarts. These settings are stored with the results.
Matching function evaluations
equalises objective-call budgets, not wall-clock time or the suitability of a parameter setting
to each function.

Table~\ref{tab:t2-cec} reports means, sample standard deviations, best and worst final objective
values. Convergence plots show median best-so-far error with an interquartile band, and
Fig.~\ref{fig:t2-cecdist} displays all final-run errors. We also apply two-sided Mann--Whitney
tests with Holm correction across the 12 pairwise comparisons. These test distributional
differences; they are not tests of mean performance, and median differences alone do not
fully describe skewed distributions.

\begin{table*}[t]
\caption{Official CEC2005 instances: final objective values over 15 runs per cell. Lower is better. Standard deviation uses $n-1$. Optima are 390 (F6) and $-330$ (F9).}
\label{tab:t2-cec}\centering\small
\begin{tabular}{@{}llrrrrr@{}}\toprule
Function & $D$ & Algorithm & Mean & Std. dev. & Best & Worst\\\midrule
F6 & 2 & GA & 396.728071 & 20.776039 & 390.000003 & 469.881943\\
F6 & 2 & PSO & 390.000000 & 0.000000 & 390.000000 & 390.000000\\
F6 & 2 & SA & 396.144878 & 6.539262 & 390.305091 & 410.602213\\
\midrule
F6 & 10 & GA & 465.166890 & 86.859403 & 395.757002 & 715.564526\\
F6 & 10 & PSO & $4.031\times10^{5}$ & $8.337\times10^{5}$ & 390.237147 & $2.014\times10^{6}$\\
F6 & 10 & SA & 1506.071423 & 2048.856085 & 428.638385 & 7670.107816\\
\midrule
F9 & 2 & GA & -330.000000 & 0.000000 & -330.000000 & -330.000000\\
F9 & 2 & PSO & -330.000000 & 0.000000 & -330.000000 & -330.000000\\
F9 & 2 & SA & -329.976837 & 0.047067 & -329.999955 & -329.844730\\
\midrule
F9 & 10 & GA & -329.992401 & 0.005209 & -329.997721 & -329.981570\\
F9 & 10 & PSO & -325.031999 & 5.041317 & -329.005041 & -311.751448\\
F9 & 10 & SA & -310.443145 & 10.497144 & -322.738317 & -288.878826\\
\bottomrule\end{tabular}\end{table*}


\subsection{Results and interpretation}
On \textbf{F6 at $D=2$}, PSO reaches the optimum to numerical precision and has the smallest
mean error. GA's mean value is 396.728, slightly worse than SA's 396.145, because one poor
GA run raises its mean. Yet GA's median is 390.025 versus SA's 394.243, and the distributions
differ after correction ($p_{\mathrm{Holm}}=0.0045$). Mean and typical-run comparisons therefore
answer different questions even in two dimensions. PSO differs from both alternatives after
correction ($p_{\mathrm{Holm}}<0.001$).

On \textbf{F9 at $D=2$}, both GA and PSO reach values extremely close to $-330$. Their
corrected comparison does not meet the 0.05 threshold ($p_{\mathrm{Holm}}=0.0559$), and their
practical difference is negligible at the displayed precision. Both improve substantially on
SA's typical final error, although SA also gets close to the optimum in some runs.

On \textbf{F6 at $D=10$}, the ranking is particularly sensitive to summary choice. GA has the
lowest mean objective, 465.167, while PSO's mean is approximately $4.03\times10^5$. However,
PSO's median is 391.661, lower than GA's 468.858. Poor PSO runs dominate its mean and spread.
The GA--PSO distributional comparison is not significant after correcting all 12 tests
($p_{\mathrm{Holm}}=0.0559$). This does not establish equivalence: it shows why a claim that
GA is unconditionally superior would exceed the evidence. SA's median is 750.341, and its
distribution differs from both GA and PSO after correction.

On \textbf{F9 at $D=10$}, GA is best on both mean and median error: mean error above the
optimum is 0.00760, compared with 4.968 for PSO and 19.557 for SA. All three pairwise
comparisons remain significant after correction. The advantage is clearer here than on
Rosenbrock because it is not explained by a few extreme failures in one competitor.

\begin{figure*}[t]\centering
\includegraphics[width=.85\textwidth]{fig14_convergence_cec2005.png}
\caption{Official CEC2005 instances: median best-so-far error and interquartile band over
15 runs. Each cell uses the same evaluation budget across algorithms. Errors below $10^{-12}$
are displayed at that floor for the logarithmic axis.}
\label{fig:t2-cecconv}\end{figure*}

\begin{figure*}[t]\centering
\includegraphics[width=.98\textwidth]{fig18_benchmark_distributions.png}
\caption{All 180 final-run errors with per-cell medians. The large PSO tail on F6 at $D=10$
explains the disagreement between mean and median rankings. A $10^{-12}$ plotting floor is
used for errors at numerical precision.}
\label{fig:t2-cecdist}\end{figure*}

These observations concern the implemented algorithms and fixed parameter settings. Tuning
the annealing schedule, changing PSO's velocity cap or adding restarts, or changing GA variation could alter the
outcomes and would constitute a new experiment. The benchmark does not reproduce the
membership-function constraints or inference cost of the controller; its rankings cannot simply
be transferred to the assistive-care tuning problem. Runtime also differs across algorithms,
so equal evaluation budgets should be interpreted as search-effort control, not equal elapsed time.

\section{Conclusions}
\label{sec:t2-conclusion}
The Mamdani controller provides a demonstrable mapping from room conditions and user
preferences to heating/cooling and lighting actions. Membership functions, a linguistic rule
base, rule-activation examples and response surfaces make the design inspectable. The
reference controller and MATLAB implementation support reproducible examination of the
chosen scenario; simulation alone does not establish effectiveness for real residents.

Genetic tuning reduces normalised test error by 32.6\% relative to the hand-designed baseline
on the synthetic reference-policy dataset. It also exposes a weakness of error-only optimisation:
three of five tuned controllers reorder linguistic sets, four leave some inputs without a defined
output, and none preserves both audited properties. The lower-error controller should therefore
be treated as an experimental result rather than a deployment-ready improvement. A stronger
next design would enforce coverage and ordering as feasibility constraints, then compare error
among valid controllers using the same training/test split.

The official CEC benchmark comparison reinforces the need to inspect distributions. PSO is
strong in two dimensions; GA gives the lowest mean objective in both ten-dimensional cells,
but the F6 median and tail tell a different story. Corrected rank tests qualify claims that
would appear decisive from means alone. No universal winner follows from two functions,
two dimensions and 15 repetitions.

The central limitations are the synthetic target policy, absence of real-resident evaluation,
small number of controller-tuning runs and fixed optimiser parameters. The work demonstrates
the requested design, optimisation and comparison processes while making the consequences
of their assumptions visible.

\section*{Individual Contributions}
Bibek Paudyal (250288) coordinated the fuzzy-controller design and Python implementation.
Sajan Mahat (250289) coordinated genetic tuning, the matched control and the coverage and
linguistic-ordering audits. Siddhartha Bhatta (250620) coordinated the MATLAB implementation
and benchmark comparison. All three members contributed equally to report preparation and
review of the implementation, experiments and conclusions, with shared responsibility for
the final submission.

